\section{Conclusion}

In this work, we addressed the challenge of long-term memory in MBRL. While prior approaches have primarily focused on improving world model architectures to better exploit contextual information, we instead explored an orthogonal direction: encoding memory directly into model parameters. To this end, we introduced TT-MBRL, a framework that leverages gradient-based updates at test time to incorporate episode-specific information into the world model. By enabling the agent to learn a distinct set of parameters for each episode, TT-MBRL moves beyond shared representations and allows rapid adaptation to the current episode, effectively forming a form of neural memory. Through experiments on both a simplified setting and the Memory Maze benchmark, we demonstrated that TT-MBRL consistently improves performance across multiple world model architectures. Notably, our method achieves state-of-the-art results on Memory Maze. We find that the effectiveness of TT-MBRL scales proportionally to the number of environment steps collected at test-time. These results highlight the potential of parameter-based memory as a powerful complement to contextual memory in MBRL systems. More broadly, our findings suggest that test-time learning constitutes a promising direction for developing agents capable of long-term planning.

\subsection{Limitations and Future Work.}

The primary limitation of our approach is the reduced evaluation efficiency resulting from the use of gradient-based optimization at test time. TT-MBRL requires learning a distinct world model and behavior policy for each evaluation episode, leading to increased computational cost and longer evaluation times. While this overhead remains manageable in environments such as Memory Maze, it may become prohibitive in settings that require a large number of environment interactions. As a direction for future work, we plan to extend our approach to operate during the exploration phase. We hypothesize that integrating TT-MBRL into exploration could improve sample efficiency and overall performance by enabling the collection of more informative and diverse trajectories. Consequently, improving the computational efficiency of our approach will be crucial to making it practical for online adaptation during exploration.